\documentclass[11pt]{article}

\usepackage[margin=1in]{geometry}
\usepackage{booktabs}
\usepackage{array}
\usepackage{longtable}
\usepackage{hyperref}
\usepackage{xcolor}
\usepackage{enumitem}
\usepackage{amsmath}
\usepackage[T1]{fontenc}
\usepackage{lmodern}

\hypersetup{
  colorlinks=true,
  linkcolor=blue!55!black,
  urlcolor=blue!55!black
}

\setlist[itemize]{leftmargin=1.5em}
\setlist[enumerate]{leftmargin=1.7em}
\renewcommand{\arraystretch}{1.18}

\title{Daily Experimental Findings: Frame-Invariant Forecasting with GRPO}
\author{InvariantForecasting Project}
\date{April 30, 2026}

\begin{document}
\maketitle

\begin{abstract}
This note summarizes the experimental findings from the April 30, 2026
RunPod session for the InvariantForecasting project. The core question was
whether GRPO-style post-training on ForecastBench paraphrase groups can reduce
forecast instability across semantically equivalent question variants while
preserving ordinary forecasting accuracy. The strongest clean result observed
today was a Mantic-style GRPO run with training invariance weight
$\lambda_{\mathrm{train}}=1$, selected by validation checkpoint performance and
evaluated once on the held-out test split. It reduced test paraphrase dispersion
from 0.08434 to 0.08083 and slightly reduced FRIE-1 from 0.32031 to 0.31938,
while increasing Brier score from 0.23597 to 0.23854. A stronger
$\lambda_{\mathrm{train}}=2$ run looked better on validation at checkpoint 10
but did not transfer to test, where paraphrase dispersion and FRIE-1 worsened.
The main conclusion is that the implementation is now much safer and the
research signal is plausible but still fragile. The next scientific step is
replication across seeds and a prespecified checkpoint-selection rule.
\end{abstract}

\section{Purpose of the Session}

The session had three practical goals:

\begin{enumerate}
  \item Verify that the GRPO implementation can train and evaluate without
  collapsing into degenerate completions such as repeated punctuation.
  \item Compare the base Qwen2.5-7B-Instruct forecaster against
  invariance-regularized GRPO variants.
  \item Determine whether the observed paraphrase-invariance improvements are
  strong enough to justify a workshop-paper direction.
\end{enumerate}

The most important methodological constraint was to distinguish a theoretical
failure from an implementation failure. In other words, if the model fails to
improve, that failure should be attributable to the training objective or data
regime, not to a broken reward parser, unstable RL configuration, or invalid
checkpoint.

\section{Dataset and Evaluation Setup}

The data are ForecastBench current-event questions filtered to questions
resolved after Qwen's relevant knowledge cutoff, with five paraphrase variants
per underlying question group.

\begin{center}
\begin{tabular}{lrrr}
\toprule
Split & Groups & Variants & Notes \\
\midrule
Train & 478 & 2,390 & Used for GRPO post-training \\
Validation & 60 & 300 & Used for checkpoint selection \\
Test & 60 & 300 & Used for final held-out comparison \\
\bottomrule
\end{tabular}
\end{center}

Each group contains five paraphrased variants of the same forecasting question.
The evaluation reports both per-variant forecast accuracy and within-group
paraphrase instability.

\section{Metrics}

The key metrics used today were:

\begin{itemize}
  \item \textbf{Brier score}: lower is better. This measures forecast accuracy
  against the realized binary outcome.
  \item \textbf{ParaSD}: lower is better. This measures dispersion of predicted
  probabilities across paraphrases in the same question group.
  \item \textbf{FRIE-1}: lower is better. This is the combined metric with
  invariance penalty weight 1, conceptually
  \[
    \mathrm{FRIE}_1 = \mathrm{Brier} + \mathrm{ParaSD}.
  \]
  \item \textbf{Coverage}: should be 1.0. This confirms all variants/groups were
  parseable and included in the metric computation.
\end{itemize}

Because the test split has only 60 groups, small absolute differences should be
treated as directional evidence until replicated across seeds.

\section{Implementation Findings}

\subsection{Original GRPO Collapse}

The original full lambda0 run was not merely a weak model. It was invalid. The
model learned to output repeated exclamation marks, yielding completions like:

\begin{quote}
\texttt{!!!!!!!!!!!!!!!!!!!!!!!!!!!!!!!!!!!!!!!!!!!!!!!!}
\end{quote}

The diagnostic logs showed the failure mode clearly:

\begin{itemize}
  \item reward parse rate dropped from 1.0 to 0.2 and then to 0;
  \item gradients became non-finite or effectively unusable;
  \item reward became a constant parse-failure value;
  \item KL and entropy became unavailable or meaningless;
  \item all later evaluations from those checkpoints had zero parseable groups.
\end{itemize}

Therefore, the original full lambda0 checkpoints should not be used for any
scientific claim.

\subsection{Safety Improvements}

The safer GRPO path introduced or relied on the following checks:

\begin{itemize}
  \item parse-rate monitoring during training;
  \item punctuation-loop detection;
  \item non-finite gradient checks;
  \item conservative learning rate, KL regularization, low sampling temperature,
  and gradient clipping;
  \item preflight validation of split sizes, prompt rendering, token lengths,
  and GRPO group layout.
\end{itemize}

The later Mantic-style runs completed with full parseability:

\begin{itemize}
  \item \texttt{variant\_coverage = 1.0}
  \item \texttt{n\_parseable\_groups = 60}
  \item \texttt{n\_variant\_predictions = 300}
\end{itemize}

This substantially increases confidence that the remaining results reflect the
objective and data regime rather than an obvious implementation crash.

\section{Main Results}

\subsection{Held-Out Test Comparison}

The most important comparison is the held-out test split. The base model was
plain Qwen2.5-7B-Instruct without GRPO post-training. Lower is better for Brier,
ParaSD, and FRIE-1.

\begin{center}
\begin{tabular}{lrrrr}
\toprule
Run & Brier $\downarrow$ & ParaSD $\downarrow$ & FRIE-1 $\downarrow$ & Coverage \\
\midrule
Base Qwen2.5-7B-Instruct & 0.235966 & 0.084341 & 0.320307 & 1.0 \\
Mantic GRPO $\lambda_{\mathrm{train}}=1$, ckpt20 & 0.238545 & 0.080833 & 0.319378 & 1.0 \\
Mantic GRPO $\lambda_{\mathrm{train}}=2$, ckpt10 & 0.235530 & 0.085824 & 0.321354 & 1.0 \\
\bottomrule
\end{tabular}
\end{center}

\subsection{Test Delta Relative to Base}

\begin{center}
\begin{tabular}{lrrr}
\toprule
Run & $\Delta$Brier & $\Delta$ParaSD & $\Delta$FRIE-1 \\
\midrule
Mantic GRPO $\lambda_{\mathrm{train}}=1$, ckpt20 & +0.002578 & -0.003507 & -0.000929 \\
Mantic GRPO $\lambda_{\mathrm{train}}=2$, ckpt10 & -0.000437 & +0.001483 & +0.001047 \\
\bottomrule
\end{tabular}
\end{center}

The lambda1 run produced the most relevant positive result: a meaningful
reduction in paraphrase dispersion and a small improvement in FRIE-1, at the
cost of worse Brier score. The lambda2 run slightly improved Brier but worsened
both paraphrase dispersion and FRIE-1 on test.

\section{Validation Checkpoint Selection}

Checkpoint selection was performed on validation, not on test. This matters
because RL checkpoints are not monotonic: later checkpoints are more trained,
but not necessarily better. The validation split was used to select a checkpoint
before test evaluation.

\subsection{Mantic GRPO, $\lambda_{\mathrm{train}}=1$}

\begin{center}
\begin{tabular}{lrrr}
\toprule
Checkpoint & Brier $\downarrow$ & ParaSD $\downarrow$ & FRIE-1 $\downarrow$ \\
\midrule
Base validation & 0.240697 & 0.084712 & 0.325409 \\
ckpt5 & 0.240153 & 0.085576 & 0.325729 \\
ckpt10 & 0.239263 & 0.080787 & 0.320050 \\
ckpt15 & 0.236637 & 0.091164 & 0.327800 \\
ckpt20 & 0.234193 & 0.084778 & 0.318971 \\
ckpt25 & 0.235847 & 0.092573 & 0.328419 \\
ckpt30 & 0.235615 & 0.085031 & 0.320646 \\
\bottomrule
\end{tabular}
\end{center}

The validation-selected checkpoint by FRIE-1 was ckpt20. On test, ckpt20 reduced
ParaSD and slightly improved FRIE-1, but worsened Brier.

\subsection{Mantic GRPO, $\lambda_{\mathrm{train}}=2$}

\begin{center}
\begin{tabular}{lrrr}
\toprule
Checkpoint & Brier $\downarrow$ & ParaSD $\downarrow$ & FRIE-1 $\downarrow$ \\
\midrule
Base validation & 0.240697 & 0.084712 & 0.325409 \\
ckpt5 & 0.237883 & 0.086303 & 0.324187 \\
ckpt10 & 0.236280 & 0.080537 & 0.316817 \\
ckpt15 & 0.235993 & 0.088916 & 0.324909 \\
ckpt20 & 0.239687 & 0.091482 & 0.331169 \\
ckpt25 & 0.238705 & 0.094720 & 0.333425 \\
ckpt30 & 0.241588 & 0.080897 & 0.322485 \\
\bottomrule
\end{tabular}
\end{center}

The validation-selected checkpoint by FRIE-1 was ckpt10. However, ckpt10 did not
transfer positively to test: test Brier improved slightly, but ParaSD and FRIE-1
worsened. This is evidence that the stronger invariance weight may overfit the
validation paraphrase structure or perturb calibration in a way that does not
generalize.

\section{Interpretation}

\subsection{What Was Discovered}

The session produced four concrete findings.

\begin{enumerate}
  \item \textbf{The original collapse was an implementation/training failure.}
  The repeated-exclamation behavior was not a meaningful learned forecasting
  policy. It was a degenerate RL failure mode caused by parse failures and
  unstable reward dynamics.

  \item \textbf{The safer GRPO implementation can run cleanly.} Later
  Mantic-style runs reached full evaluation coverage with parseable outputs and
  no observed punctuation-loop collapse.

  \item \textbf{A lambda1 invariance objective can reduce test paraphrase
  sensitivity.} The best clean test result reduced ParaSD from 0.08434 to
  0.08083. This is the strongest positive empirical signal from the day.

  \item \textbf{More invariance pressure is not automatically better.} Lambda2
  looked excellent on validation but failed to improve test FRIE-1. This is a
  useful negative result and should shape the next experimental design.
\end{enumerate}

\subsection{Scientific Strength of the Result}

The result is promising but not yet conclusive. The lambda1 test improvement in
ParaSD is aligned with the research hypothesis, and the fact that the evaluation
remained fully parseable makes it usable. However:

\begin{itemize}
  \item the test split contains only 60 groups;
  \item seed variance is visible in earlier lambda1 replications;
  \item Brier sometimes worsens when ParaSD improves;
  \item validation improvements do not always transfer to test;
  \item no confidence intervals or bootstrap tests have been computed yet.
\end{itemize}

The appropriate claim at this point is therefore:

\begin{quote}
In a controlled ForecastBench paraphrase setting, invariance-regularized GRPO
can reduce paraphrase sensitivity in some runs while preserving parseability,
but the effect is modest and seed/checkpoint sensitive.
\end{quote}

This is a plausible workshop-paper seed, not yet a completed ICML-level result.

\section{Recommended Next Steps}

\subsection{Immediate Experimental Steps}

\begin{enumerate}
  \item Preserve the current RunPod artifacts: \texttt{results/}, \texttt{logs/},
  \texttt{configs/}, and \texttt{outputs/best/}.
  \item Run bootstrap confidence intervals over groups for Brier, ParaSD, and
  FRIE-1.
  \item Replicate the Mantic lambda1 and lambda2 setups over multiple seeds.
  \item Select checkpoints by validation FRIE-1 only, then evaluate exactly once
  on test.
  \item Inspect qualitative examples where ParaSD improved and where it worsened.
\end{enumerate}

\subsection{Paper-Framing Steps}

If replications preserve the direction of the lambda1 result, a credible
workshop paper could focus on:

\begin{itemize}
  \item paraphrase instability as a measurable forecasting reliability problem;
  \item ForecastBench paraphrase groups as a compact evaluation bed;
  \item GRPO with an invariance penalty as a lightweight intervention;
  \item the observed tradeoff between calibration/accuracy and paraphrase
  invariance;
  \item negative results showing that stronger invariance penalties can fail to
  generalize.
\end{itemize}

\section{Bottom Line}

Today did discover something real. The most important result is not that the
method is already solved, but that the project now has:

\begin{itemize}
  \item a known invalid failure mode and safeguards against it;
  \item a clean evaluation pipeline with full parseability;
  \item a positive lambda1 test signal on paraphrase dispersion;
  \item a negative lambda2 generalization result that constrains the theory.
\end{itemize}

The project has enough signal to justify continued experimentation and possibly
the beginning of a workshop-paper outline, provided the next round focuses on
replication, uncertainty estimates, and careful checkpoint-selection protocol.

\end{document}
